% !TeX root = ../nc-quadric-classification.tex
% 2026-11-21 / v1.32（証明本文をBへ移管）
\section{相対 cubic tails の導来認識と効果性}
\label{sec:relative-cubic-tails}

\subsection{三次解消による相対 tails の compact 認識}

以下では $R$ を任意の可換環とする。$R$ の正則性も
Noetherian 性も、まず仮定しない。
$C$ は正の向きの $R$-線形 $\mathbb Z$-代数、すなわち
$C_{ij}=0$（$j<i$）、$C_{ii}=R$ であり、積を
$C_{ij}\otimes_R C_{jk}\to C_{ik}$ と書くものとする。
右加群 $P_i=e_iC$ の成分は $(P_i)_j=C_{ij}$ である。
$S_i$ は頂点 $i$ にだけ $R$ を置く右加群とする。
次数 torsion は、各元が十分遠い頂点への全ての射で
消されるという意味で用いる。

\begin{theorem}[相対 cubic tails の四項 tilting]
\label{thm:relative-cubic-tails-tilting}
$C$ は隣接成分で生成され、全ての $C_{ij}$ は
有限 projective $R$-加群であるとする。
有限 projective $R$-加群 $V_i,R_i$ と可逆 $R$-加群 $W_i$ により
\[
 0\longrightarrow W_i\otimes_RP_{i+4}
 \longrightarrow R_i\otimes_RP_{i+3}
 \longrightarrow V_i\otimes_RP_{i+1}
 \longrightarrow P_i\longrightarrow S_i\longrightarrow0
 \tag{CT-res}\label{eq:cubic-tails-simple-resolution}
\]
が全ての $i$ で完全であると仮定する。さらに
\[
 \operatorname{Ext}^q_{\operatorname{Gr}C}(S_l,P_j)=0
 \quad\text{unless }q=3\text{ and }j=l+4
 \tag{CT-dual}\label{eq:cubic-tails-relative-dual}
\]
とし、$\mathcal L_l=\operatorname{Ext}^3(S_l,P_{l+4})$ は
可逆 $R$-加群であるとする。
$\mathcal O_i=\pi P_i\in\operatorname{QGr}C$ と置く。
このとき次が成立する。
\begin{enumerate}
\item $\mathbf R\omega$ の cohomological dimension は高々二であり、
  $D(\operatorname{QGr}C)$ は $\mathcal O_i$（$i\in\mathbb Z$）で
  compactly generated される。
\item 任意の $i,j$ について
  \[
  \operatorname{Ext}^q_{\operatorname{QGr}C}
       (\mathcal O_i,\mathcal O_j)
   \simeq
   \begin{cases}
     C_{ji},&q=0,\\
     \operatorname{Hom}_R(C_{i,j-4},\mathcal L_{j-4}),&q=2,\ j\ge i+4,\\
     0,&\text{otherwise}.
   \end{cases}
  \tag{CT-coh}\label{eq:cubic-tails-line-cohomology}
  \]
  とくに任意の連続四項の直和
  $T_a=\bigoplus_{i=a}^{a+3}\mathcal O_i$ は compact tilting generator
  である。その通常の自己準同型環を $\Gamma_a$ とすると
  \[
    D(\operatorname{QGr}C)\simeq D(\operatorname{Mod}\Gamma_a),\qquad
    D(\operatorname{QGr}C)^c\simeq\operatorname{Perf}\Gamma_a.
    \tag{CT-Morita}\label{eq:cubic-tails-canonical-morita}
  \]
\end{enumerate}
全ての仮定と有限 projective 係数の解消が任意の $R$-基底変更で
保たれる場合、これらの compact 比較もその基底変更と整合する。
本定理は $C$ の Noetherian 性や coherent heart の存在を結論しない。
\end{theorem}
\begin{proof}
証明は相対モジュライ論文 \cite[\texttt{relative_cubic_tails}]{QuadricModuli2026} に置く。
\end{proof}

\begin{lemma}[有限型 heart の一様な cohomology]
\label{lem:relative-cubic-coherent-cohomology}
前定理で $R$ を Noetherian とする。
さらに $\operatorname{QGr}C$ が locally Noetherian であり、
$\mathcal O_i$ が Noetherian な生成系をなすと仮定する。
任意の Noetherian 対象 $M$ について
\[
 \operatorname{Ext}^q(\mathcal O_i,M)
 \text{ は有限 }R\text{-加群},\qquad
 \operatorname{Ext}^{q}(\mathcal O_i,M)=0
 \quad(q>0,\ i\gg0).
 \tag{CT-finite}\label{eq:cubic-tails-coherent-finiteness}
\]
任意の Noetherian 対象 $M,N$ の Hom も有限 $R$-加群である。
同じ仮定を満たす平坦基底変更の下では、これらの Hom は基底変更と可換する。
cofinal な偶数添字部分列による生成条件も成立するなら、
その二項移動は Artin--Zhang ample である。
\end{lemma}
\begin{proof}
証明は相対モジュライ論文 \cite[\texttt{relative_cubic_tails}]{QuadricModuli2026} に置く。
\end{proof}

\subsection{右軌道規約と既存 foundation への適用}

ここで既存の $D_{ij}=\operatorname{Hom}(L_i,L_j)$ と右軌道環
$B$ の規約を明示的に合わせる。正の向きの右加群に使う
$\mathbb Z$-代数は
\[
 C_{ij}=D_{-j,-i}=\operatorname{Hom}(L_{-j},L_{-i})\quad(i\le j)
 \tag{CT-right}\label{eq:cubic-tails-right-convention}
\]
であり、積は射の合成 $f\otimes g\mapsto f\circ g$ である。
これは $D$ の反対代数を添字反転したもの。
もとの $D$ の右 tails と、添字を変えず同一視してはならない。
この規約では $\pi P_i$ が $L_{-i}$ に対応し、
$i=-3,-2,-1,0$ の四項窓が元の foundation を与える。

\begin{proposition}[幾何的 cubic 開の canonical Grothendieck 導来比較]
\label{prop:quadric-canonical-cubic-tails}
命題~\ref{prop:quadric-uniform-cubic-hom-open} の有限型開を用いる。
任意の affine 基底変更 $T=\operatorname{Spec}R\to U^\circ$ について
自然な導来 Morita 比較
\[
 D(\operatorname{QGr}B_T)\simeq D(\operatorname{Mod}\Lambda_T),\qquad
 D(\operatorname{QGr}B_T)^c\simeq\operatorname{Perf}\Lambda_T
 \tag{QT-compact}\label{eq:quadric-tails-compact-comparison}
\]
を得る。この比較は構造対象を $F_0$ に、次数移動 $(1)$ を
既存の $\mathsf s$ に対応させ、任意の導来基底変更と整合する。
特に compact 部分は相対 smooth proper であり、
その canonical Serre 関手について
$(1)^2\simeq S^{-1}[2]$ が成立する。
これは Noetherian 性を仮定しない Grothendieck tails の結論である。
Noetherian 基底では、後の定理~\ref{thm:quadric-finite-type-effectivity} が
$B_T$ の左右 Noetherian 性を証明するので、
$\operatorname{QGr}B_T=\operatorname{Ind}(\operatorname{qgr}B_T)$
も成立する。
\end{proposition}
\begin{proof}
証明は相対モジュライ論文 \cite[\texttt{relative_cubic_tails}]{QuadricModuli2026} に置く。
\end{proof}

\begin{corollary}[Noetherian な心上の Hom と豊富性]
\label{cor:quadric-noetherian-base-recognition}
上の affine 基底 $R$ を Noetherian とし、$B_T$ を左右 Noetherian
と仮定する。このとき
$\operatorname{QGr}B_T=\operatorname{Ind}(\operatorname{qgr}B_T)$ であり、
前命題はこの heart に固有の canonical compact 導来圏の比較となる。
一般の coherent 対象間の Hom は有限で、平坦基底変更と可換する。
$\pi B_T$ は基底上平坦かつ exceptional、$(\pi B_T,(1))$ は AZ ample である。
これらの主張は左右 Noetherian 性が保たれる全有限型基底変更で成立する。
特に系~\ref{cor:quadric-point-local-noetherian-as4} の
任意の Noetherian 局所基底 $R$ と、その任意の有限型可換 $R$-代数で
成立する。後の定理~\ref{thm:quadric-finite-type-effectivity} により、
単一の有限型開全体と任意の Noetherian 基底変更にも適用できる。
\end{corollary}
\begin{proof}
証明は相対モジュライ論文 \cite[\texttt{relative_cubic_tails}]{QuadricModuli2026} に置く。
\end{proof}

\begin{lemma}[相対 tails の一様 cohomology から全ての $\chi_q$ へ]
\label{lem:relative-cubic-global-chi}
系~\ref{cor:quadric-noetherian-base-recognition} の Noetherian 基底と
右 Noetherian 環 $B=B_T$ において、任意の有限生成次数付き右加群 $M$ と
任意の $q\ge0$ に対して
\[
   \underline{\operatorname{Ext}}^q_B(R,M)
   \quad\text{は有限 }R\text{-加群である}.
\]
従って右相対 $\chi_q$ が全ての $q$ で成立する。
反対環にも同じ cubic 議論を適用すれば左側も成立する。
この結論は、左右 Noetherian 性がある同じ affine 基底全体での主張である。
\end{lemma}
\begin{proof}
証明は相対モジュライ論文 \cite[\texttt{relative_cubic_tails}]{QuadricModuli2026} に置く。
\end{proof}

\subsection{相対内在的条件への昇格}

\begin{theorem}[Noetherian な基底上での相対条件の全確認]
\label{thm:quadric-noetherian-relative-promotion}
命題~\ref{prop:quadric-canonical-cubic-tails} の幾何的 cubic 開
$U^\circ$ への affine 射 $\operatorname{Spec}R\to U^\circ$ を取り、
$R$ は Noetherian 可換環、$B_R$ は左右 Noetherian 環であると仮定する。
任意の可換有限型 $R$-代数 $R'$ に対して
\[
 \mathcal C_{R'}=\operatorname{qgr}(B_R\otimes_RR'),
 \qquad\mathcal O_{R'}=\pi(B_R\otimes_RR')
\]
と置く。この系は、有限型 $R$-基底の site 上で
定義~\ref{def:relative-abelian-system} および
定義~\ref{def:relative-h2} の全ての条件を満たす。
その局所的な平方根は次数移動 $(1)$ であり、
canonical compact 導来圏は
\[
 D(\operatorname{Ind}\mathcal C_{R'})^c
     \simeq\operatorname{Perf}(\Lambda_R\otimes_RR')
\]
である。この比較は module-object 引戻しの導来関手による
任意の有限型基底変更と整合する。

従って $\operatorname{Spec}R$ 自身が元の局所有限型 $k$-site に属する場合、
$(\mathcal C_R,\mathcal O_R)$ は実際に
$\mathfrak H_2(R)$ の対象を与える。
また系~\ref{cor:quadric-point-local-noetherian-as4} の任意の
Noetherian 局所基底 $R$ は本定理の仮定を満たす。
その $R$ が元の site に属さない場合にも、有限型 $R$-基底の系に関する
上の全条件は成立する。ただしこの記述だけで
$\mathfrak H_2$ の一般基底への極限保存的な延長を定義したとはしない。
\end{theorem}
\begin{proof}
証明は相対モジュライ論文 \cite[\texttt{relative_cubic_tails}]{QuadricModuli2026} に置く。
\end{proof}
